% ============================================================================
% FULL CORRECTED \section{Methodology} for the major revision.
% Diff against the original. Changes reflect the actual leak-free runs:
%  - cohort N~118, ~23 positive (~20%), NOT balanced 63/63
%  - EEG = frozen LaBraM foundation-model embeddings (200-D), full montage;
%    hand-crafted (Hjorth + spectral entropy) kept only as a baseline
%  - eye 11-D, behaviour 6-D; d_model = 128 (compact, ~2.9M params)
%  - repeated subject-level CV (5x10), ROC-AUC primary (imbalance)
%  - foundation-model EEG training regimes (linear probe / window-level / FT)
%  - external validation on HBN; calibration/referral; recent baselines
% New cite keys to add to .bib: LaBraM2024, HBNEEG2024, DrFuse2024,
%   ConformalTriage2025.
% ============================================================================
\section{Methodology}
\label{method}
\subsection{Dataset and Experimental Protocol}
\subsubsection{Dataset}
This study used the Multimodal Resource for Studying Information Processing in the
Developing Brain (MIPDB)~\cite{Langer2017}, a publicly available multimodal dataset
released by the Child Mind Institute (CMI). MIPDB provides EEG, eye-tracking, and
behavioural task-performance recordings collected across a six-paradigm cognitive
battery, alongside clinician-administered DSM-5 diagnostic assessments. We extract
each modality stream from the paradigm in which it is most informative and integrate
them at the participant level (Section~\ref{sec:tasks}), rather than at the level of
trial-locked synchronized events.

The analysis cohort comprised $N \approx 118$ participants (aged 6--44 years) with a
usable EEG recording and a diagnostic label; of these, 23 met criteria for at least
one DSM-5 diagnosis (approximately 20\% prevalence) and the remainder had no recorded
diagnosis. The complete-case multimodal subset, with all three modalities present,
comprised $\approx 104$ participants. The cohort size is modest relative to
large-scale clinical studies and reflects the practical difficulty of acquiring
multimodal neurophysiological data; the class imbalance is characteristic of
screening settings. Because the cohort is imbalanced, we use ROC-AUC as the primary
metric and report balanced accuracy as a secondary measure
(Section~\ref{sec:partition}). The diagnosis group was heterogeneous (ADHD, Autism
Spectrum Disorder, anxiety, and depressive disorders); individual disorder labels
were not used during training. The classification target is the binary distinction
between any confirmed diagnosis and the absence of any diagnosis. This formulation
tests whether multimodal markers generalize across heterogeneous clinical
presentations rather than detecting a single disorder, consistent with the Research
Domain Criteria framework~\cite{Insel2014}; we frame NeuroFormer as a
screening-support tool rather than a standalone diagnostic system.

\textbf{Participant selection.} From MIPDB participants with diagnostic information,
we retained those with (i) a usable EEG recording passing quality control
(Section~\ref{sec:eegpipe}) and (ii) a valid diagnostic label. Eye-tracking and
behavioural features were additionally extracted where available; participants
missing a non-EEG modality are retained via the missing-modality regime
(Section~\ref{sec:missing}) rather than discarded, which mitigates the
complete-case selection bias of the original design.

\subsubsection{Experimental Tasks and Modality Sourcing}
\label{sec:tasks}
The MIPDB battery contains six paradigms: Resting State, Sequence Learning, Symbol
Search, Surround Suppression, Contrast Change Detection (a two-alternative
forced-choice task, denoted SAIIT 2AFC), and Naturalistic Stimuli (four short video
clips)~\cite{Langer2017}. Because not every paradigm yields informative signal in
every modality, we sourced each modality from the paradigm(s) where its signal is
richest:
\begin{itemize}
    \item \textbf{EEG:} aggregated across all six paradigms within each participant,
    giving a participant-level summary of cortical activity over the session.
    \item \textbf{Eye-tracking:} extracted from the Naturalistic Stimuli paradigm,
    where sustained viewing produces rich oculomotor and pupil variability.
    \item \textbf{Behavioural:} extracted from the Contrast Change Detection (SAIIT
    2AFC) paradigm, which yields trial-level reaction-time and accuracy data.
\end{itemize}
This produces three participant-level feature streams per individual. Fusion in
NeuroFormer therefore operates at the \emph{participant} level, not at
trial-synchronized events; Section~\ref{conclusion} discusses this design choice.

\subsubsection{Data Partitioning and Cross-Validation}
\label{sec:partition}
All results use \textbf{strict subject-level} partitioning: every recording from a
participant is confined to a single fold, eliminating subject-level leakage. We use
\textbf{repeated stratified $k$-fold cross-validation} (5 folds $\times$ 10 repeats),
preserving the class prevalence in each fold. We report ROC-AUC as the primary metric
(appropriate to the $\sim$20\% prevalence) with 95\% confidence intervals across
folds, and use \emph{paired} comparisons (identical splits) across feature sets and
models. All fold-dependent steps (standardization, calibration, any feature fitting)
are fit on training folds only. Given the cohort size we deliberately avoid heavy
hyperparameter search, which would invite optimistic bias.

To assess generalization beyond MIPDB, we externally validate on the independent
\textbf{Healthy Brain Network (HBN)} EEG cohort~\cite{HBNEEG2024}
($N \approx 1{,}980$, RestingState), harmonized to the same preprocessing
(Section~\ref{sec:external}).

This study used a publicly available, fully anonymized dataset collected under
institutional ethical approval at the Child Mind Institute
(\url{http://fcon_1000.projects.nitrc.org/indi/cmi/eeg/}). We release preprocessing
code, extracted features, and trained checkpoints on GitHub.

\subsection{Modality-Specific Feature Extraction}
\begin{figure*}[t]
\centering
\includegraphics[width=0.95\textwidth]{figures/feature_extraction_pipelines.png}
\caption{Modality-specific preprocessing and feature-extraction pipelines. EEG (top)
is aggregated across the six MIPDB paradigms within each participant and mapped to
\textbf{frozen foundation-model embeddings}; eye-tracking features (middle) are
extracted from the Naturalistic Stimuli paradigm and behavioural features (bottom)
from the Contrast Change Detection paradigm. The streams produce participant-level
representations of dimensionality \textbf{200} (LaBraM EEG embedding), \textbf{11}
(eye), and \textbf{6} (behaviour), linearly projected to a shared model dimension
$d_{\text{model}} = 128$ before the transformer (Section~\ref{method}).}
\label{fig:pipeline}
\end{figure*}
Each modality is processed in non-overlapping windows over its paradigm-specific
recording, producing per-window feature vectors that are linearly projected to the
model dimension $d_{\text{model}} = 128$ and augmented with sinusoidal positional
encoding. The result is three participant-level sequences
$\mathbf{E} \in \mathbb{R}^{T_{\text{EEG}} \times 128}$,
$\mathbf{Y} \in \mathbb{R}^{T_{\text{Eye}} \times 128}$, and
$\mathbf{Z} \in \mathbb{R}^{T_{\text{Beh}} \times 128}$ (Fig.~\ref{fig:pipeline}).
Because each sequence comes from a different paradigm, the temporal index $t$ does
not align to a shared physical event across modalities; cross-modal attention in
NeuroFormer captures participant-level relationships rather than trial-locked
dependencies.

\subsubsection{EEG Pipeline}
\label{sec:eegpipe}
MIPDB EEG was acquired at 500\,Hz with a 128-channel EGI HydroCel Geodesic Sensor
Net~\cite{Langer2017}. We retain the \textbf{full montage} (vertex reference renamed
$\mathrm{E129}\!\to\!\mathrm{Cz}$, GSN-HydroCel-129 positions attached for spatial
interpolation), concatenate all six paradigms within each participant, band-pass
filter to 0.1--75\,Hz, apply a 60\,Hz notch, resample to 200\,Hz, and scale to
microvolts. Signals are segmented into non-overlapping \textbf{15\,s windows}
(3000 samples) and clipped at $\pm 800\,\mu$V to bound gross artifacts. A
quality-control gate retains only physiological recordings (windowed median
amplitude in 3--150\,$\mu$V, $<$15\% clipped), discarding flat, saturated, and
impedance/calibration files.

Each 15\,s window is mapped to a \textbf{200-dimensional embedding} by a
\textbf{frozen, pretrained EEG foundation model (LaBraM)}~\cite{LaBraM2024}, using
the \texttt{braindecode} \texttt{InterpolatedLaBraM} implementation, whose
interpolation front-end projects the recorded montage onto LaBraM's canonical
channels. Embeddings are extracted with no gradient updates to the backbone. For
fixed-vector analyses, per-window embeddings are aggregated per participant by
\textbf{mean$+$std pooling} (concatenated mean and standard deviation across
windows), which outperformed mean-only pooling.

\textbf{Hand-crafted EEG baseline.} For a controlled comparison we also compute a
classical descriptor on the same windows---Hjorth activity, mobility, and complexity
together with Welch spectral entropy per channel---representing the
pre-foundation-model approach.

\subsubsection{Foundation-Model EEG and Training Regimes}
\label{sec:regimes}
Under identical subject-level CV we compare how the frozen LaBraM representation is
\emph{used}: (i) a \textbf{subject-pooled linear probe} on the mean$+$std embedding;
(ii) \textbf{window-level training} of a classification head on the frozen backbone,
in which each 15\,s window is a training instance and per-participant predictions are
the mean window probability (multiple-instance learning); and (iii) \textbf{end-to-end
fine-tuning} of the full backbone. Window-level training on the frozen model is the
recommended configuration; end-to-end fine-tuning overfits at this cohort size
(Section~\ref{results}).

\subsubsection{Eye-Tracking Pipeline}
Binocular gaze and pupil signals were recorded at 120\,Hz with an SMI iView~X RED-m
eye-tracker, 5-point calibrated before each paradigm~\cite{Langer2017}. From the
Naturalistic Stimuli paradigm we compute, per block, an \textbf{11-dimensional}
feature vector: pupil diameter (mean and standard deviation, left/right), gaze
position (mean and standard deviation of point-of-regard $x$/$y$), head-position
variability ($x$/$y$/$z$), response latency (mean and standard deviation), and the
gaze vector (mean and standard deviation of its components). Block-level features are
pooled (mean) per participant. Values were validated against the original pipeline
(per-feature $r \approx 1.0$).

\subsubsection{Behavioural Pipeline}
From the Contrast Change Detection paradigm (SAIIT 2AFC) we compute, per participant,
\textbf{6} summary statistics over valid trials: mean, standard deviation, median,
minimum, and maximum reaction time, and overall accuracy~\cite{Macmillan2004}. The
6-dimensional vector is linearly projected to $d_{\text{model}} = 128$ and combined
with a learned modality-type embedding before the shared encoder.

We use feature-level rather than raw end-to-end representations for the non-EEG
modalities because the cohort size makes raw sequence modelling prone to overfitting
at transformer scale, and because compact descriptors are robust to the
cross-paradigm sourcing of the three streams. For EEG, the foundation-model
embedding provides the learned representation that the original hand-crafted
features lacked.

\subsection{Model Architecture}
\begin{figure*}[t]
\centering
\includegraphics[width=0.95\textwidth]{figures/neuroformer_architecture_diagram.png}
\caption{NeuroFormer architecture (compact variant). The EEG stream enters as frozen
LaBraM embeddings; each modality stream receives a learned modality-type token
(\texttt{[EEG]}, \texttt{[Eye]}, \texttt{[Beh]}) and sinusoidal positional encoding,
and passes through modality-specific 2-layer transformer encoders ($H{=}8$,
$d_{\text{model}}{=}128$). The encoded sequences are concatenated with a learnable
\texttt{[CLS]} token and processed by a 2-layer cross-modal fusion transformer; the
fused \texttt{[CLS]} representation passes through two fully connected layers
(FC$_1$: $128{\to}512$; FC$_2$: $512{\to}256$, ReLU, dropout $p{=}0.3$) and a softmax
output. Modality dropout enables any-subset inference.}
\label{fig:NeuroFormer}
\end{figure*}

NeuroFormer processes the three feature streams with a two-stage transformer
architecture (Fig.~\ref{fig:NeuroFormer}). Modality-specific encoders learn
within-modality representations; a cross-modal fusion transformer then models
inter-modality dependencies and produces the prediction. Lightweight reconstruction
decoders are attached to the encoder outputs as a self-supervised regularizer. The
contribution lies in adapting hierarchical multimodal fusion---fed by
foundation-model EEG and trained under a missing-modality regime---to a constrained
clinical setting, rather than proposing a new transformer primitive. We report a
\textbf{compact} variant ($\approx$2.9\,M parameters) to match the cohort size and
directly address the overparameterization concern.

\subsubsection{Modality Encoders}
The representations $\mathbf{E}$, $\mathbf{Y}$, $\mathbf{Z}$ enter
\textbf{modality-specific} encoders (separate parameters per modality) with their
positional and modality-type embeddings applied. Each is a 2-layer transformer
encoder with 8 attention heads and feed-forward dimension $d_{\text{ff}} = 256$
\cite{Vaswani2017}, producing encoded sequences
$\tilde{\mathbf{E}}, \tilde{\mathbf{Y}}, \tilde{\mathbf{Z}}$.

\subsubsection{Cross-Modal Fusion Transformer}
The three encoded sequences are concatenated along the sequence dimension and a
learnable \texttt{[CLS]} token~\cite{Devlin2019} is prepended, giving length
$1 + T_{\text{EEG}} + T_{\text{Eye}} + T_{\text{Beh}}$. The modality-type embeddings
(\texttt{[EEG]}, \texttt{[Eye]}, \texttt{[Beh]}) mark each token's modality, whereas
\texttt{[CLS]} aggregates across all modalities for the prediction. A 2-layer,
8-head transformer lets \texttt{[CLS]} attend across the three streams.

\subsubsection{Missing-Modality Training}
\label{sec:missing}
Because modalities are frequently absent in real screening, we train with
\textbf{modality dropout}: each modality is randomly withheld per sample (never all),
and the modality-type embeddings act as availability indicators, so one model serves
any subset of $\{$EEG, eye, behaviour$\}$ with graceful degradation, following
missing-modality clinical-fusion work~\cite{DrFuse2024}.

\subsubsection{Classification Head and Interpretability}
The fused \texttt{[CLS]} representation $\mathbf{f}_{\text{CLS}} \in \mathbb{R}^{128}$
passes through two fully connected layers (FC$_1$: $128 \to 512$, FC$_2$:
$512 \to 256$) with ReLU and dropout ($p = 0.3$), then a softmax output. For
interpretability we aggregate the final-fusion-layer attention from \texttt{[CLS]} to
each modality's tokens, yielding participant-level modality-importance scores. As the
streams lack a shared physical time origin, we report attention as a participant-level
summary of which streams drive each prediction, treated as a post-hoc explanatory aid
rather than mechanistic evidence (Section~\ref{conclusion}).

\subsubsection{Architectural Hyperparameters}
The compact variant uses $d_{\text{model}} = 128$, $H = 8$ heads ($d_k = 16$),
$d_{\text{ff}} = 256$, $N_{\text{enc}} = 2$ modality-encoder layers, and
$N_{\text{fus}} = 2$ fusion layers; dropout $p = 0.3$ throughout. Values were chosen
via the architectural-depth ablations in Section~\ref{ablation}, where reduced
capacity matched or exceeded the original 18.9\,M-parameter configuration at this
cohort size.

\subsubsection{Training and Optimization}
For the multimodal model we minimise a composite objective combining binary
cross-entropy with auxiliary modality-specific reconstruction:
\begin{equation}
\mathcal{L} \;=\; \mathcal{L}_{\text{CE}} \;+\; \lambda \!\!\sum_{m \in \{\text{EEG},\,\text{Eye},\,\text{Beh}\}}\!\! \mathcal{L}_{\text{rec}}^{(m)},
\end{equation}
where $\mathcal{L}_{\text{CE}}$ is the class-weighted binary cross-entropy and
$\mathcal{L}_{\text{rec}}^{(m)}$ the mean-squared reconstruction error from a
lightweight decoder at each encoder output; $\lambda$ is selected on validation
folds. The reconstruction term regularizes the encoders toward structured signal
content~\cite{Vincent2008}; we report its ablation, including on the external set
(Section~\ref{results}). For the EEG-only window-level model
(Section~\ref{sec:regimes}) we train the classification head on frozen LaBraM
embeddings with class-weighted cross-entropy and mixed precision.

Parameters are optimized with Adam~\cite{Kingma2015} (initial learning rate
$10^{-4}$, weight decay $10^{-4}$). The multimodal variant uses cosine-annealed warm
restarts~\cite{Loshchilov2017}, dropout ($p = 0.3$), early stopping, and feature
augmentation. The compact model contains $\approx$2.9\,M parameters. We implement the
pipeline in PyTorch~2.11 with \texttt{braindecode}~$\geq$1.5 and MNE-Python~1.12 on
the Digital Research Alliance of Canada (Nibi) cluster (NVIDIA H100 for fine-tuning;
CPU for frozen extraction and the compact fusion model), with fixed random seeds for
reproducibility.

\subsection{Calibrated, Referral-Aware Prediction}
To support triage, out-of-fold scores are calibrated post-hoc with temperature
scaling; we report the Expected Calibration Error before and after. We add
split-conformal prediction and a risk--coverage analysis: low-confidence cases
abstain and are referred for clinician review, and accuracy/AUC on the retained
(confident) subset is reported as a function of coverage, following conformal
clinical-triage methods~\cite{ConformalTriage2025}.

\subsection{Baselines}
Under the identical subject-level protocol we compare: single-modality classifiers,
naive concatenation, stacking late fusion, the NeuroFormer attention fusion, a
\textbf{single-stage multimodal Transformer} (modality-as-token; the no-hierarchy
control), and a \textbf{gated attention fusion}. The frozen LaBraM front-end itself
constitutes the self-supervised / foundation-model baseline. All baselines use the
same features, splits, and metric, with paired significance testing.

\subsection{External Validation on HBN}
\label{sec:external}
Frozen LaBraM embeddings are extracted from HBN RestingState with the same
preprocessing and quality control. As HBN provides continuous CBCL psychopathology
factors rather than categorical diagnoses, we predict each factor from the
embeddings (Ridge regression for Pearson $r$; logistic regression on the median split
for ROC-AUC) under the same subject-level CV, and additionally apply the
MIPDB-trained model directly (transfer). All evaluations are at the well-powered
scale $N \approx 1{,}980$, and we report the out-of-distribution result together with
the label-construct and paradigm differences from MIPDB.
